% ============================================================
% Section: Related Work
% ============================================================

\section{Related Work}
\label{sec:related}

% TODO: Write related work (~1 page, 3–4 subsections)

% ---- 2.1 -----------------------------------------------
\subsection{Blind Image Quality Assessment}
\label{sec:related:iqa}

No-reference image quality assessment (NR-IQA) methods predict perceptual quality without access to a clean reference image.
BRISQUE~\cite{mittal2012making} extracts hand-crafted natural scene statistics and outputs a single scalar quality score, making it unsuitable for conditioning downstream models on fine-grained quality signals.
NIMA~\cite{talebi2018nima} learns to predict the distribution of human aesthetic ratings but still collapses quality into one dimension and is trained on general photographic datasets with no exposure to clinical imaging artifacts.
More recent CLIP-based approaches such as CLIP-IQA leverage vision-language alignment for zero-shot quality estimation, yet their general-domain training yields poor calibration on dermoscopic imaging where domain-specific distortions (contact artifacts, uneven illumination, specular highlights) dominate.
In contrast, VisiScore-Net decomposes quality into five clinically motivated dimensions—focus, illumination, contrast, completeness, and artifact—providing a differentiable, multi-dimensional quality vector that downstream modules can directly condition on.

% ---- 2.2 -----------------------------------------------
\subsection{Uncertainty Estimation in Deep Learning}
\label{sec:related:uncertainty}

Principled uncertainty quantification is essential for safety-critical deployment of medical AI.
MC Dropout~\cite{gal2016dropout} approximates Bayesian inference by sampling stochastic forward passes at test time, but its uncertainty estimates are not informed by input quality, causing overconfident predictions on severely degraded images.
Deep Ensembles~\cite{lakshminarayanan2017simple} train multiple independent networks and aggregate their predictive distributions; while empirically strong, this approach is computationally expensive and likewise provides no mechanism to modulate uncertainty based on observed image quality.
The Variational Information Bottleneck (VIB)~\cite{alemi2017deep} offers a principled information-theoretic compression framework: the encoder is trained to retain only task-relevant features by penalizing the KL divergence between the latent posterior and a fixed Gaussian prior.
However, a fixed prior cannot reflect the varying reliability of the input signal—a low-quality image should induce a broader, higher-entropy posterior than a high-quality one.
Our Quality-conditioned VIB (Q-VIB) replaces the fixed prior with a quality-adaptive prior whose variance is a learned function of the VisiScore vector, directly embedding quality information into the KL regularization; we prove (Proposition~2) and empirically verify on three datasets that predictive entropy decreases monotonically with quality ($\rho = -0.165$, $p < 10^{-120}$ on ISIC).

% ---- 2.3 -----------------------------------------------
\subsection{Image Restoration and Enhancement}
\label{sec:related:enhancement}

Learning-based image restoration has advanced rapidly.
NAFNet~\cite{chen2022simple} simplifies non-linear activation functions for efficient denoising and deblurring, achieving state-of-the-art PSNR on standard benchmarks.
Restormer~\cite{zamir2022restormer} introduces a transformer-based architecture with multi-scale attention, excelling at deraining and deblurring tasks.
Diffusion-based methods such as DiffBIR further improve perceptual quality by leveraging generative priors.
Despite strong distortion metrics, none of these methods constrain the restoration process to preserve task-relevant features: improving PSNR/SSIM does not guarantee that the enhanced image retains the diagnostic information critical for downstream classification, and the change in AUC ($|\Delta\mathrm{AUC}|$) is left unconstrained.
VisiEnhance-Net addresses this gap with a Diagnosis-Preserving Loss (DP-Loss) that aligns the latent representations of enhanced and reference images in the Q-VIB feature space, providing a principled bound on diagnostic drift.

% ---- 2.4 -----------------------------------------------
\subsection{Dermoscopic Image Analysis and Triage}
\label{sec:related:dermoscopy}

Automated dermoscopy analysis has matured considerably, with deep neural networks achieving or exceeding dermatologist-level performance on curated benchmarks such as ISIC~2018 and HAM10000.
Yet top-performing systems implicitly assume that input images meet a minimum quality threshold; in practice, images captured by non-specialist users with smartphone dermoscopes exhibit substantial quality variation, and AUC degrades markedly on low-quality subsets.
Existing triage workflows treat quality control as a hard binary gate—either the image is accepted or the user is asked to retake it—with no intermediate repair pathway.
Our dual-channel agent closes this gap by routing images through a quality-conditioned pipeline: recoverable degradations are corrected by VisiEnhance-Net before diagnosis, while unrecoverable images (flagged by low completeness score $q_3$) trigger a retake request, explicitly modeling the quality-diagnosis relationship rather than assuming it away.
